% Vol III --- Gluing Chapter, Part IV: Van den Bergh derived Morita gluing.
% Architecture: notes/gluing_chapter_architecture_2026_04_23.md, Part IV (§§4.1--4.4).
% Primary sources:
%   Van den Bergh 2004 (arXiv:math/0211064), Non-commutative crepant resolutions
%   Bondal--Kapranov 1990 (Math. USSR Izv.), Enhanced triangulated categories
%   Bondal--Orlov 2002 (arXiv:math/0206295), Derived categories of coherent sheaves
%   King 1997, Moduli of representations of finite-dimensional algebras
%   Davison--Hennecart--Schlegel Mejia 2023 (arXiv:2303.12592), BPS Lie algebras
%   Davison 2017 (arXiv:1601.02479), Critical CoHA and PBW
%   Derksen--Weyman--Zelevinsky 2008 (arXiv:0704.0649), Quivers with potentials
%   Keller--Yang 2011 (arXiv:0906.0761), Derived equivalences from mutations
%   Szendroi 2008 (arXiv:0705.3419), Non-commutative DT on resolved conifold
%   Craw--Ishii 2004, and Craw 2008 on (-3)-curve tilting
%   Kontsevich--Soibelman 2008 (arXiv:0811.2435), Motivic DT and wall-crossing
%   Bridgeland 2006 (arXiv:math/0502050), Flops and derived equivalence
%   Ginzburg 2006 (arXiv:math/0612139), Calabi--Yau algebras

\providecommand{\C}{\mathbb{C}}
\providecommand{\Z}{\mathbb{Z}}
\providecommand{\N}{\mathbb{N}}
\providecommand{\Q}{\mathbb{Q}}
\providecommand{\R}{\mathbb{R}}
\providecommand{\bP}{\mathbb{P}}
\providecommand{\cO}{\mathcal{O}}
\providecommand{\cA}{\mathcal{A}}
\providecommand{\cC}{\mathcal{C}}
\providecommand{\cD}{\mathcal{D}}
\providecommand{\cE}{\mathcal{E}}
\providecommand{\cF}{\mathcal{F}}
\providecommand{\cG}{\mathcal{G}}
\providecommand{\cH}{\mathcal{H}}
\providecommand{\cM}{\mathcal{M}}
\providecommand{\cT}{\mathcal{T}}
\providecommand{\cU}{\mathcal{U}}
\providecommand{\cY}{\mathcal{Y}}
\providecommand{\Tot}{\mathrm{Tot}}
\providecommand{\Coh}{\mathrm{Coh}}
\providecommand{\Perf}{\mathrm{Perf}}
\providecommand{\Hom}{\mathrm{Hom}}
\providecommand{\End}{\mathrm{End}}
\providecommand{\Ext}{\mathrm{Ext}}
\providecommand{\RHom}{R\mathrm{Hom}}
\providecommand{\mod}{\mathrm{mod}}
\providecommand{\Mod}{\mathrm{Mod}}
\providecommand{\Spec}{\mathrm{Spec}}
\providecommand{\Proj}{\mathrm{Proj}}
\providecommand{\Jac}{\mathrm{Jac}}
\providecommand{\Crit}{\mathrm{Crit}}
\providecommand{\tr}{\mathrm{tr}}
\providecommand{\id}{\mathrm{id}}
\providecommand{\vphi}{\varphi}
\providecommand{\gl}{\mathfrak{gl}}
\providecommand{\fg}{\mathfrak{g}}
\providecommand{\ghat}{\widehat{\mathfrak{gl}}_{1}}
\providecommand{\Yplus}{Y^{+}}
\providecommand{\CoHA}{\mathrm{CoHA}}
\providecommand{\Hilb}{\mathrm{Hilb}}
\providecommand{\NCCR}{\mathrm{NCCR}}
\providecommand{\Mu}{\mathrm{M}}
\providecommand{\kch}{\kappa_{\mathrm{ch}}}
\providecommand{\kcat}{\kappa_{\mathrm{cat}}}
\providecommand{\kBKM}{\kappa_{\mathrm{BKM}}}
\providecommand{\kfib}{\kappa_{\mathrm{fibre}}}
\providecommand{\ClaimStatusTheorem}{\textsc{[theorem]}}
\providecommand{\ClaimStatusDefinition}{\textsc{[definition]}}
\providecommand{\ClaimStatusDerivation}{\textsc{[first-principles]}}
\providecommand{\ClaimStatusConjectured}{\textsc{[conjectural]}}
\providecommand{\ClaimStatusProposition}{\textsc{[proposition]}}

% =======================================================================
\part{Van den Bergh derived Morita gluing}
\label{part:gluing-vdb}
% =======================================================================

Parts I--III assembled global CoHAs from \v{C}ech, toric, and McKay atlases: each
overlap carried a geometric transition --- a Fourier--Mukai kernel, a toric
cocycle, a BKR equivalence. Part IV replaces the atlas by a single object. A
\emph{tilting bundle} $T$ on a projective CY$_3$ (or a local CY$_3$)
collapses the derived category of coherent sheaves onto the derived category of
modules over a finite-dimensional algebra $A = \End(T)$: all gluing
information on $X$ is encoded in the multiplicative structure of $A$. The Hall
product on critical cohomology transports along this collapse. Mutations of $T$
act as the chamber-changing automorphisms of the associated CoHA.

% =======================================================================
\section{Van den Bergh tilting bundles}
\label{sec:vdb-tilting}
% =======================================================================

\subsection*{The obstruction: $D^b(\Coh X)$ is not compactly generated from below}

For $X$ a smooth projective variety, $D^b(\Coh X)$ admits no canonical
generator. Locally one trivialises $\Coh X$ by $\Coh U \simeq \mod R_U$ with
$R_U = \Gamma(U,\cO_X)$ a commutative affine ring, but the local presentations
are incompatible: $R_U$ is Morita-trivial (its module category is generated by
the regular module $R_U$ itself), whereas $D^b(\Coh X)$ globally admits no such
single-generator description. A tilting bundle is the object that repairs this
failure \emph{globally}.

\begin{definition}[Tilting bundle; Bondal--Kapranov~1990]
\label{def:tilting-bundle}
\ClaimStatusDefinition\ Let $X$ be a smooth variety. A coherent sheaf
$T \in \Coh(X)$ is a \emph{tilting bundle} if
\begin{enumerate}
\item[(T1)] $T$ is locally free of finite rank;
\item[(T2)] $\Ext^{i}_{\cO_X}(T, T) = 0$ for all $i \geq 1$
           (``vanishing higher self-Ext'');
\item[(T3)] $T$ generates $D^{b}(\Coh X)$ as a triangulated category:
           the smallest full triangulated subcategory of $D^{b}(\Coh X)$ closed
           under shifts, cones, and direct summands and containing $T$ is
           $D^{b}(\Coh X)$ itself.
\end{enumerate}
The algebra $A = \End_{\cO_X}(T)^{\mathrm{op}}$ is called the \emph{tilting
algebra}.
\end{definition}

Conditions (T1)--(T2) force $\RHom(T, T) = \End(T) = A$ concentrated in degree
zero; condition (T3) is the generation property. The reason for passing to the
opposite algebra is that $\RHom(T, -)$ produces a \emph{right} $\End(T)$-module
and, by convention, one wishes to write these as left $A$-modules.

\begin{theorem}[Derived Morita via tilting; Bondal 1989; Rickard 1989 for algebras;
Bondal--Van den Bergh 2003, arXiv:\texttt{math/0204218} Thm.~3.1.1]
\label{thm:derived-morita-tilting}
\ClaimStatusTheorem\ Let $T$ be a tilting bundle on $X$ and $A = \End(T)^{\mathrm{op}}$.
The pair
\[
  L \;=\; T \otimes^{L}_{A} (-) \;:\; D(A\text{-}\mod) \;\longrightarrow\; D_{\mathrm{qc}}(X),
  \qquad
  R \;=\; \RHom_{\cO_X}(T, -) \;:\; D_{\mathrm{qc}}(X) \;\longrightarrow\; D(A\text{-}\mod)
\]
is an adjoint pair $L \dashv R$. Two hypotheses make it an equivalence after
restriction to the compact part:
\begin{enumerate}
\item[(H1)] (\emph{Compact generation}) $T$ is a compact generator of
$D_{\mathrm{qc}}(X)$: the smallest localising subcategory of $D_{\mathrm{qc}}(X)$
closed under arbitrary coproducts and containing $T$ is all of $D_{\mathrm{qc}}(X)$,
and $T$ is compact (Bondal--Van den Bergh \emph{loc.\ cit.}, Thm.~3.1.1);
\item[(H2)] (\emph{Higher Ext vanishing}) $\Ext^{>0}_{\cO_X}(T, T) = 0$,
so that the unit $A \to R(L(A)) = \RHom(T, T \otimes^L_A A) = \RHom(T, T)$ is an
isomorphism in $D(A\text{-}\mod)$ concentrated in degree zero.
\end{enumerate}
Under (H1)--(H2), $L \dashv R$ is a pair of mutually quasi-inverse equivalences
\[
  D_{\mathrm{qc}}(X) \;\simeq\; D(A\text{-}\mod), \qquad
  \Perf(X) \;\simeq\; \Perf(A) = D^{\mathrm{perf}}(A\text{-}\mod).
\]
Restricting to bounded coherent objects on the left and finitely presented
modules on the right yields the classical equivalence
$D^{b}(\Coh X) \simeq D^{b}(A\text{-}\mod)$.
\end{theorem}

\begin{proof}[Proof.]
The adjunction $L \dashv R$ is the standard tensor--Hom adjunction for the
$A$-$\cO_X$ bimodule $T$: for $M \in D(A\text{-}\mod)$ and $\cF \in D_{\mathrm{qc}}(X)$,
$\Hom_{D_{\mathrm{qc}}(X)}(T \otimes^L_A M, \cF) = \Hom_{D(A)}(M, \RHom_{\cO_X}(T, \cF))$
functorially. The unit $\eta_A \colon A \to R(L(A)) = \RHom(T, T)$ is an
isomorphism by (H2), which forces $\Ext^{>0}(T, T) = 0$ and identifies $\RHom(T,T)
\simeq \End(T) = A^{\mathrm{op}}$ concentrated in degree zero; since $A$ generates
$D(A\text{-}\mod)$, this upgrades to $\eta_M \colon M \xrightarrow{\sim} R(L(M))$
for all $M \in D(A\text{-}\mod)$. The counit $\varepsilon_\cF \colon L(R(\cF)) \to
\cF$ is an isomorphism on the compact generator $T$ by the same computation, and
$T$ compactly generates $D_{\mathrm{qc}}(X)$ by (H1), so $\varepsilon$ is an
isomorphism on all of $D_{\mathrm{qc}}(X)$ (Bondal--Van den Bergh
arXiv:\texttt{math/0204218}, Thm.~3.1.1: a compact generator with bounded Ext
produces an equivalence with modules over its endomorphism dg-algebra, here
concentrated in degree zero by (H2)).
\end{proof}

\subsection*{Non-commutative crepant resolutions}

The obstruction on singular $X$ is sharper: $X$ admits no smooth model within
its birational class without blowing up, but one may ask for a
\emph{non-commutative} smooth substitute.

\begin{definition}[NCCR; Van den Bergh~2004, arXiv:\texttt{math/0211064}]
\label{def:nccr}
\ClaimStatusDefinition\ Let $R$ be a commutative Noetherian normal Gorenstein
domain. A \emph{non-commutative crepant resolution} of $R$ is an $R$-algebra
$A = \End_R(M)$, for some reflexive $R$-module $M$, such that
\begin{enumerate}
\item[(N1)] $A$ is \emph{homologically homogeneous}: all simple $A$-modules have
the same projective dimension, equal to $\dim R$;
\item[(N2)] $A$ is a \emph{maximal Cohen--Macaulay} $R$-module;
\item[(N3)] $A$ is generically Azumaya and $M$ is a reflexive generator.
\end{enumerate}
\end{definition}

``Crepant'' means $\omega_A \cong A$ as $A$-bimodules: the canonical class is
trivial, matching the CY condition on $\Spec R$.

\begin{theorem}[Van den Bergh equivalence, arXiv:\texttt{math/0211064} Thm.~6.6.1
and \S8]
\label{thm:vdb-equivalence}
\ClaimStatusTheorem\ Let $Y \to X = \Spec R$ be a crepant resolution (in the
commutative sense) with $R$ as in Definition~\ref{def:nccr}, and suppose
$\dim R \leq 3$. Then there exists a tilting bundle $T_Y$ on $Y$ with
$\End(T_Y)^{\mathrm{op}} = A$ an NCCR of $R$ in the sense of
Definition~\ref{def:nccr}, and the adjoint pair $L \dashv R$ of
Theorem~\ref{thm:derived-morita-tilting} is an equivalence
\[
  D^{b}(\Coh Y) \;\simeq\; D^{b}(A\text{-}\mod).
\]
Both sides realise the same abstract triangulated category; when $R$ is
Gorenstein of dimension three with trivial dualising module, the
\emph{Ginzburg dg-lift} $\Gamma = \Gamma(Q, W)$ of $A$ (Ginzburg 2006,
arXiv:\texttt{math/0612139} Thm.~3.2.8) carries a CY$_{3}$ structure:
$\Gamma$ is exact $3$-Calabi--Yau as a dg-algebra, meaning the inverse
dualising complex satisfies
$\RHom_{\Gamma^e}(\Gamma, \Gamma \otimes \Gamma)[3] \simeq \Gamma$ in the
derived category of $\Gamma$-bimodules.

\emph{Scope of existence.} The hypothesis $\dim R \leq 3$ is sharp. Beyond
dimension three, the existence of an NCCR is not guaranteed and \emph{fails}
in the cases most of interest: projective K3 surfaces with generic polarisation
(where $H^{0,2}(X) \neq 0$ obstructs higher Ext-vanishing), the quintic and
other generic projective CY$_3$ without toric or orbifold presentation, and
compact hyperk\"ahler manifolds of dimension $\geq 4$. The Van den Bergh
existence theorem covers precisely the local three-dimensional crepant setting:
toric CY$_3$, orbifolds $\C^n/G$ with $G \subset SU(n)$ of dimension $n \leq 3$,
$(-3)$-curve contractions, and their relatives.
\end{theorem}

On the conifold $R = \C[x,y,z,w]/(xw - yz)$ the NCCR is the \emph{Klebanov--Witten
quiver algebra} with tilting bundle $T = \cO \oplus \cO(1)$ on the resolved
conifold $\tilde Y = \cO_{\bP^1}(-1) \oplus \cO_{\bP^1}(-1)$. On local $\bP^2 =
\Tot(\cO_{\bP^2}(-3))$ the NCCR is the nine-arrow Beilinson algebra, with
$T = \cO \oplus \cO(1) \oplus \cO(2)$. On $(-3)$-curve contractions
$\tilde Y \to X$ with $\tilde Y$ containing a $\bP^1$ of self-intersection
$(-1,-1,-1,-1)$ on its normal bundle, Craw's tilting bundles (Craw
2008, building on Craw--Ishii 2004) play the same role. In each case, an
explicit tilting object $T$ reduces all geometric CoHA questions on $X$ to
finite-dimensional Hall-algebra questions on $A = \End(T)^{\mathrm{op}}$; the
superpotential $W$ reconstructing $A$ as a Jacobi algebra is uniquely
determined (up to right-equivalence) from the cyclic $A_\infty$-structure on
the tilting algebra (Bocklandt 2008, arXiv:\texttt{math/0603558} Thm.~3.1).

\subsection*{Existence and failure}

The Van den Bergh theorem (Theorem~\ref{thm:vdb-equivalence}) guarantees the
existence of a tilting bundle in precisely the \emph{local three-dimensional}
crepant setting: $X$ obtained as a crepant resolution of an affine Gorenstein
singularity $\Spec R$ with $\dim R \leq 3$. This covers toric CY$_3$,
orbifolds $\C^n/G$ with $G \subset SU(n)$ of dimension $n \leq 3$,
$(-3)$-curve contractions, and their relatives. Beyond this scope, global
tilting \emph{fails}. Three cases are central:
\begin{itemize}
\item[(a)] \emph{Projective K3 surface $S$ with generic polarisation.} The
obstruction is $\Ext^{2}_{\cO_S}(\cF, \cF) \neq 0$ for any locally-free $\cF$:
Serre duality on a CY$_2$ identifies
$\Ext^2(\cF, \cF) \cong \Hom(\cF, \cF)^{\vee} \neq 0$. The higher-Ext vanishing
hypothesis (H2) of Theorem~\ref{thm:derived-morita-tilting} cannot hold. The
K3 category admits compact generation (by Bondal--Van den Bergh, Thm.~3.1.1)
but no tilting generator.
\item[(b)] \emph{Quintic threefold $X \subset \bP^4$.} The obstruction is the
non-triviality of $H^{0,3}(X) = \C$ (the holomorphic volume form), which
combined with Serre duality forces $\Ext^3_{\cO_X}(\cF, \cF) \neq 0$ for
non-zero $\cF$. No $\Ext^{>0}$-vanishing tilting generator exists globally.
\item[(c)] \emph{Compact hyperk\"ahler manifold of dimension $\geq 4$.}
Hyperk\"ahler symmetry forces $h^{0,p}(X) = 1$ for all $p$ even, $0 \leq p \leq
\dim X$; the higher $\Ext$-vanishing hypothesis is obstructed at every even
degree $\leq \dim X$.
\end{itemize}
The general pattern: at $d \geq 3$, if $H^{0,q}(X) \neq 0$ for some $1 \leq q
\leq d$, then $\Ext^q(T, T) \supseteq H^{0,q}(X) \neq 0$ for any
locally-free $T$ with $\cO_X$ as a summand; the hypothesis (H2) of
Theorem~\ref{thm:derived-morita-tilting} fails. In such cases, no global
tilting exists; only local tilting on Zariski patches, and the gluing of
Parts I--III (or Part V via factorisation algebras) applies.

% =======================================================================
\section{Tilting as gluing presentation}
\label{sec:tilting-as-gluing}
% =======================================================================

The tilting bundle is a \emph{presentation} of $X$ by a single
quiver-with-potential: the quiver vertices are the summands of
$T = \bigoplus_i T_i$, the arrows are extensions among them, and the potential
encodes the exact $3$-Calabi--Yau constraint on the Ginzburg dg-lift
$\Gamma(Q, W)$.

\subsection*{The decomposition $T = \bigoplus_i T_i$}

Write $T = \bigoplus_{i \in I} T_i$ as a finite direct sum of indecomposable
summands. This decomposition is canonical up to isomorphism of summands
(Krull--Schmidt on the additive category generated by $T$). Then
\[
  A \;=\; \End(T)^{\mathrm{op}} \;=\; \bigoplus_{i, j \in I}
  \Hom_{\cO_X}(T_i, T_j)^{\mathrm{op}}
\]
is a matrix algebra over the ``diagonal'' algebras $\End(T_i)$, indexed by
pairs $(i,j)$. The off-diagonal block $\Hom_{\cO_X}(T_i, T_j)$ is the
\emph{arrow space} from $i$ to $j$: its basis is the space of morphisms
$T_i \to T_j$ in $\Coh(X)$ modulo scaling.

In the conifold example $T = \cO \oplus \cO(1)$:
\[
  \Hom(\cO, \cO(1)) \;=\; H^{0}(\bP^1, \cO(1)) \oplus H^{0}(\bP^1, \cO(-1))
  \;=\; \C^{2} \oplus 0 \;=\; \C a_1 \oplus \C a_2,
\]
and $\Hom(\cO(1), \cO) = H^0(\cO(-1)) \oplus H^0(\cO(1)) = 0 \oplus \C^{2}
= \C b_1 \oplus \C b_2$ (the second term arises via the $\cO(-1) \oplus
\cO(-1)$ fibre direction). Similarly $\End(\cO) = \End(\cO(1)) = \C$. The
Klebanov--Witten quiver with vertices $\{\circ, \bullet\}$, arrows $a_1, a_2
\colon \circ \to \bullet$ and $b_1, b_2 \colon \bullet \to \circ$, and
superpotential $W = a_1 b_1 a_2 b_2 - a_1 b_2 a_2 b_1$ emerges \emph{directly}
from the cohomology count.

\subsection*{The Jacobi algebra, the potential, and the Ginzburg dg-lift}

The endomorphism algebra of a tilting bundle on a smooth local CY$_3$ carries
a cyclic derivation structure: there exists a superpotential
$W \in \C Q_{\mathrm{cyc}} = \C Q / [\C Q, \C Q]$ in the cyclic quotient of
the path algebra such that the relations among the arrows are precisely the
cyclic derivatives $\partial_a W = 0$. The pair $(Q, W)$ is the quiver with
potential; the \emph{Jacobi algebra} is
\[
  J = \Jac(Q, W) \;=\; \C Q / \langle \partial_a W \colon a \in Q_1 \rangle.
\]
The CY$_3$ property is carried not by $J$ but by its \emph{Ginzburg dg-lift}.

\begin{definition}[Ginzburg dg-algebra; Ginzburg 2006,
arXiv:\texttt{math/0612139} \S5]
\label{def:ginzburg-dg}
\ClaimStatusDefinition\ Given a quiver with potential $(Q, W)$, the
\emph{Ginzburg dg-algebra} $\Gamma(Q, W)$ is the tensor algebra on the graded
quiver
\begin{itemize}
\item arrows of $Q$ in degree $0$;
\item one reverse arrow $a^{*} \colon j \to i$ in degree $-1$ for each
$a \colon i \to j$ in $Q_1$;
\item one loop $t_i \colon i \to i$ in degree $-2$ at each vertex $i \in I$,
\end{itemize}
equipped with the differential
$d a = 0$, $d a^{*} = \partial_a W$, and $d t_i = \sum_{a \in Q_1}
[a, a^{*}] \cdot e_i$ (the total commutator at vertex $i$).
\end{definition}

\begin{proposition}[CY-3 property lives on $\Gamma$, not on $J$;
Ginzburg 2006, Thm.~3.2.8 and \S5.3]
\label{prop:cy3-on-gamma-not-j}
\ClaimStatusProposition\ For $(Q, W)$ a quiver-with-potential arising from a
local CY$_3$:
\begin{enumerate}
\item[(a)] The dg-algebra $\Gamma(Q, W)$ is \emph{exact $3$-Calabi--Yau}: its
inverse dualising complex satisfies
$\RHom_{\Gamma^e}(\Gamma, \Gamma \otimes \Gamma)[3] \simeq \Gamma$ as
$\Gamma$-bimodules, with the CY$_3$ shift degree $+3$ matching the PTVV shift
$(d, \mathrm{shift}, E_n^{\mathrm{cl}}) = (3, -1, E_1)$ under dg-Hochschild
Serre duality.
\item[(b)] The Jacobi algebra $J = H^0(\Gamma(Q, W))$ is the degree-zero
cohomology of $\Gamma$; it is \emph{not} by itself a CY$_3$ algebra. Calling
$J$ ``CY$_3$'' is a category error: $J$ has no intrinsic shift structure. The
three-dimensional duality is an assertion about $\Gamma$-bimodules.
\item[(c)] For $Q$ connected and $W$ suitably non-degenerate,
$H^{<0}(\Gamma(Q, W)) = 0$, so $\Gamma \simeq J$ in the derived category of
$\Gamma$ (hence derived Morita between $\Gamma$ and $J$), and only via this
quasi-isomorphism does the CY$_3$ structure on $\Gamma$ descend to the
derived categorical structure $D^b(J\text{-}\mod)$.
\end{enumerate}
\end{proposition}

\begin{theorem}[Tilting-via-quiver presentation; Ginzburg 2006, Keller 2011
arXiv:\texttt{1102.4148}]
\label{thm:tilting-jacobi}
\ClaimStatusTheorem\ Let $T = \bigoplus_{i \in I} T_i$ be a tilting bundle on a
smooth local CY$_3$ $X$. Let $Q$ be the quiver with vertex set $I$ and arrow
spaces $Q_1(i,j) = \Ext^{1}_{\cO_X}(T_i, T_j)^{*}$. There exists a
quasi-homogeneous superpotential $W \in \C Q_{\mathrm{cyc}}$, unique up to
gauge, such that
\[
  \Jac(Q, W) \;=\; \End_{\cO_X}(T)^{\mathrm{op}},
  \qquad
  \Gamma(Q, W) \;\simeq\; \RHom_{\cO_X}(T, T) \text{ as dg-algebras},
\]
and the CY$_3$ structure on $\Perf(X)$ transports (via
Theorem~\ref{thm:derived-morita-tilting} applied to $\Gamma$) to the exact
CY$_3$ structure on $\Gamma$ of Proposition~\ref{prop:cy3-on-gamma-not-j}(a).
\end{theorem}

\begin{proof}[Proof sketch.]
The arrow spaces are Ext-spaces: under (H2) the underived and derived
$\Hom(T_i, T_j)$ agree in degree zero, so the first non-trivial term for
$i \neq j$ is $\Ext^{1}(T_i, T_j)$, which Serre-dualises to
$\Ext^{2}(T_j, T_i)$ on the CY$_3$; the potential encodes the $m_3$ Massey
product. The construction of $W$ is the Kontsevich--Costello
reconstruction of a cyclic $A_\infty$-algebra from its degree-zero piece plus
Serre duality: the $m_3$ on $\Ext^1$ evaluated against Serre duality
reproduces the degree-3 piece of $W$, higher $m_k$ give higher-degree
corrections, and quasi-homogeneity follows from the internal CY$_3$ grading
on $\Gamma$. The quasi-isomorphism $\RHom(T, T) \simeq \Gamma(Q, W)$ is
Keller's derived-equivalence theorem (Keller 2011 \S A.15). The CY$_3$
property, as Proposition~\ref{prop:cy3-on-gamma-not-j} records, is a property
of $\Gamma$ as a dg-algebra; the statement ``$J$ is CY$_3$'' is shorthand for
``$\Gamma$ is CY$_3$ and $J = H^0(\Gamma)$''.
\end{proof}

\subsection*{Canonical examples}

\paragraph{Conifold $\tilde Y = \cO_{\bP^1}(-1)^{\oplus 2}$.} Tilting
$T = \cO \oplus \cO(1)$. Quiver: two vertices, two arrows each way. Potential
$W = a_1 b_1 a_2 b_2 - a_1 b_2 a_2 b_1$ (Klebanov--Witten, Szendr\H{o}i 2008).

\paragraph{Local $\bP^2 = \Tot(\cO_{\bP^2}(-3))$.} Tilting
$T = \cO \oplus \cO(1) \oplus \cO(2)$ (Beilinson 1978). Quiver: three
vertices, $3 \times 3 = 9$ arrows grouped as three Koszul strands (Beilinson
triple $\cO, \cO(1), \cO(2)$). Potential
$W_{P^2} = \sum_{\sigma \in S_3} \mathrm{sgn}(\sigma)\,
x_{\sigma(1)} y_{\sigma(2)} z_{\sigma(3)}$: the antisymmetric McKay cubic
(Beasley--Plesser 2001).

\paragraph{$(-3)$-curves.} Let $\tilde Y$ contain a rational curve $C \cong
\bP^1$ with normal bundle $N_{C/\tilde Y} = \cO(-1) \oplus \cO(-1)$ that can
be contracted to a singular point of $X$. Craw 2008 (building on the
Craw--Ishii 2004 analysis of $\bC^3/G$-Hilb schemes) constructs tilting
$T_C = \cO_{\tilde Y} \oplus \cO_{\tilde Y}(C)$ (or an explicit line-bundle
refinement with intersection numbers read off from $N_{C/\tilde Y}$). The
associated tilting algebra is derived-equivalent to the quiver from the
contracted singularity; tilting automorphisms realise the flop $\tilde Y
\dashrightarrow \tilde Y^+$ as derived Morita on the same CY$_3$ category.

\paragraph{$\C^3/\Z_n$ orbifolds.} Tilting
$T = \bigoplus_{k=0}^{n-1} \cO \otimes \chi_k$ on the resolved orbifold, where
$\chi_k$ are the characters of $\Z_n$. Quiver: $n$ vertices arranged on a
cycle; arrows per McKay graph; potential cubic (type A McKay).

\subsection*{Quiver-with-potential as universal gluing data}

The force of Theorem~\ref{thm:tilting-jacobi} is that for every tilting
presentation, a single pair $(Q, W)$ plays the role of the entire \v{C}ech
atlas plus all its transition functors. One algebra $A = \Jac(Q, W)$ encodes
everything. In particular:
\begin{itemize}
\item The \v{C}ech nerve $\{U_\alpha\} \cup \{U_\alpha \cap U_\beta\} \cup \ldots$
  on $X$ collapses to the vertex set $I$ and the arrow set $Q_1$;
\item The transition functors between charts collapse to the relations
  $\partial_a W = 0$;
\item The higher cocycle data on triple overlaps collapses to the higher
  $A_\infty$ structure on $\Gamma(Q, W)$.
\end{itemize}
The tilting presentation is therefore the \emph{economy} of the gluing: one
single finite-dimensional algebra replaces the infinite-dimensional \v{C}ech
diagram.

% =======================================================================
\section{Transport of Hall product along derived Morita}
\label{sec:hall-transport}
% =======================================================================

The CoHA on $X$ is defined via the moduli stack of coherent sheaves
$\cM(X)$ with its critical cohomology ring structure (Kontsevich--Soibelman
2008, Davison--Meinhardt 2020 \texttt{arXiv:1601.02479}). On the NCCR side, the
same construction yields the CoHA on the moduli stack $\cM(A)$ of
finite-dimensional $A$-modules. The derived Morita equivalence lifts to these
moduli stacks, and Davison--Hennecart--Schlegel Mejia (DHSM) 2023
\texttt{arXiv:2303.12592} prove that the Hall product transports. The
assertion of this section is the precise form of that transport and its
first-principles derivation.

\subsection*{Moduli-stack lift of derived Morita}

\begin{theorem}[Moduli-stack Morita; Davison--Hennecart--Schlegel Mejia 2023,
arXiv:\texttt{2303.12592} Thm.~A]
\label{thm:moduli-morita}
\ClaimStatusTheorem\ Let $X$ be a smooth local CY$_3$ admitting a tilting
bundle $T$ satisfying hypotheses (H1)--(H2) of
Theorem~\ref{thm:derived-morita-tilting}, and let $A = \End(T)^{\mathrm{op}} = \Jac(Q, W)$
be the associated Jacobi algebra. Then the adjoint pair $L \dashv R$ descends
to an equivalence of (derived) moduli stacks of semistable objects
\[
  \cM^{\mathrm{ss}}(X) \;\simeq\; \cM^{\mathrm{ss}}(A)
\]
under the class-level identification
\[
  \mathrm{cl}_T \;:\; K_0(\Coh X) \;\xrightarrow{\;\sim\;}\; \Z^{I}, \qquad
  [\cF] \;\longmapsto\; (\dim \RHom(T_i, \cF))_{i \in I}.
\]
The lift intertwines the universal family $\cE_X$ on $\cM(X)$ with the
universal representation $V_A$ on $\cM(A)$: $\RHom(T, \cE_X) = V_A$ as
coherent sheaves on the common moduli stack.
\end{theorem}

\begin{proof}[Sketch.]
Derived Morita transports the abelian heart (a $t$-structure on $D^b(\Coh X)$
whose heart is the image of $A$-$\mod$) to the standard heart on $A$-$\mod$.
Bridgeland's identification of stability conditions preserves semistability
under this $t$-structure alignment. The resulting map on moduli is an
equivalence of stacks because the universal family $\cE_X$ on $\cM(X)$ maps to
the universal representation $V_A$ on $\cM(A)$, and conversely, via the
inverse functor $T \otimes^L_A -$.
\end{proof}

\subsection*{Critical cohomology transports}

The CoHA on either side is
\[
  \cH(X) \;=\; H^{\bullet,\mathrm{BM}}_{T_{\C^\times}}(\cM(X), \vphi_{\mathrm{tr}\,W_X}),
\]
where $\vphi$ is the vanishing-cycle functor and $W_X$ is the canonical CY$_3$
superpotential (pulled back from the framing). On the NCCR side, the
superpotential is precisely $W = W_{Q,W}$ --- the potential of the
quiver-with-potential --- and $\vphi_{\tr W}$ is the vanishing-cycle sheaf on
the representation variety $\mathrm{Rep}(Q, \underline{d})$.

\begin{theorem}[Vanishing-cycle transport; DHSM 2023, extending Davison 2017]
\label{thm:vanishing-cycle-transport}
\ClaimStatusTheorem\ Under the moduli equivalence of
Theorem~\ref{thm:moduli-morita}, the vanishing-cycle sheaf
$\vphi_{\tr W_X}$ on $\cM(X)$ is canonically isomorphic to
$\vphi_{\tr W_Q}$ on $\cM(A)$, compatibly with the $T$-equivariant structure.
Consequently:
\[
  \cH(X) \;\simeq\; \cH(A) \quad\text{as equivariant Borel--Moore
  cohomology modules over the respective moduli stacks.}
\]
\end{theorem}

\begin{proof}[First-principles derivation.]
A coherent sheaf $\cF$ on $X$ corresponds, under $\Phi_T$, to an $A$-module
$M = \RHom(T, \cF)$. Both $\cF$ and $M$ admit \emph{the same} deformation
theory: the obstruction space is $\Ext^{2}(\cF, \cF) \cong \Ext^{2}_A(M, M)$
(Serre duality on $X$ matches the cyclic duality on $A$). The superpotential
$W_X$ on the deformation space matches $\tr W_Q$ pulled back along the tilting
isomorphism, because
\[
  W_X|_{T^{*}_{[\cF]} \cM(X)} \;=\;
  \text{degree-3 part of the $A_\infty$ structure on $\Ext^\bullet(\cF, \cF)$},
\]
and Theorem~\ref{thm:tilting-jacobi} identifies precisely this degree-3 piece
with the Jacobi potential $W$. The vanishing-cycle sheaf is therefore an
invariant of the pair (deformation space, superpotential), transporting
identically along the moduli equivalence.
\end{proof}

\subsection*{Extension stacks and Hall product transport}

The Hall product on $\cH(X)$ is built from the \emph{extension stack}
$\mathrm{Ext}(\cM \times \cM, \cM)$: for charge classes
$\gamma_1, \gamma_2 \in K_0(\Coh X)$, the extension stack parametrises short
exact sequences $0 \to \cF_1 \to \cF \to \cF_2 \to 0$ with
$[\cF_i] = \gamma_i$ and $[\cF] = \gamma_1 + \gamma_2$. The Hall product is
convolution along the push/pull diagram
\[
  \cM_{\gamma_1} \times \cM_{\gamma_2}
  \;\xleftarrow{\;\mathrm{ends}\;}\;
  \mathrm{Ext}_{\gamma_1, \gamma_2}
  \;\xrightarrow{\;\mathrm{middle}\;}\;
  \cM_{\gamma_1 + \gamma_2}.
\]

\begin{proposition}[Extension-stack transport]
\label{prop:extension-stack-transport}
\ClaimStatusProposition\ Under the moduli equivalence
$\cM(X) \simeq \cM(A)$ the extension stacks satisfy
\[
  \mathrm{Ext}_{\gamma_1, \gamma_2}^{X} \;\simeq\;
  \mathrm{Ext}_{\mathrm{cl}_T(\gamma_1), \mathrm{cl}_T(\gamma_2)}^{A},
\]
with the ``ends'' and ``middle'' morphisms intertwined by $\Phi_T$. Combined
with Theorem~\ref{thm:vanishing-cycle-transport}, this yields:
\end{proposition}

\begin{theorem}[Hall-product transport along derived Morita; DHSM 2023
arXiv:\texttt{2303.12592} Thm.~A, with super-grading extension Thm.~B]
\label{thm:hall-product-transport}
\ClaimStatusTheorem\ Under hypotheses (H1)--(H2) on the tilting bundle $T$,
there is an isomorphism of (super-graded) bialgebras
\[
  \star_T \;:\; \cH(X) \;\xrightarrow{\;\sim\;}\; \cH(A)
\]
intertwining the Hall product, the natural coproduct (where defined), and the
$T_{\C^\times}$-equivariant structures. On charges, $\star_T$ is the
$\Z$-linear extension of $\mathrm{cl}_T$. The \emph{$\Z/2$-super grading}
(cohomological parity on $\vphi_{\tr W}$) transports along the same
equivalence, by DHSM Thm.~B.
\end{theorem}

\begin{proof}[Proof.]
The Hall product is a composition of pullback-along-ends, intersection with
the vanishing-cycle sheaf, and pushforward-along-middle. Each factor transports
under derived Morita: the ends and middle morphisms transport by
Proposition~\ref{prop:extension-stack-transport}; the vanishing-cycle sheaf
transports by Theorem~\ref{thm:vanishing-cycle-transport}; the equivariant
structure transports because the tilting bundle is $T_{\C^\times}$-equivariant
in the toric setting (and equivariantly localisable at the fixed locus in the
general local CY$_3$ setting). The $\Z/2$-super grading transports because
$\vphi_{\tr W_X}$ and $\vphi_{\tr W_Q}$ carry identically-graded sheaves of
nearby cycles under the cohomological parity; DHSM Thm.~B upgrades Thm.~A from
an isomorphism of bialgebras to an isomorphism of super-bialgebras.
\end{proof}

Theorem~\ref{thm:hall-product-transport} is established unconditionally by
DHSM (2023, Thm.~A + Thm.~B) for $A = \Jac(Q, W)$ the Jacobi algebra of any
quiver-with-potential; the CY$_3$ hypothesis is absorbed into the assertion
that the Ginzburg dg-lift $\Gamma(Q, W)$ is exact $3$-Calabi--Yau
(Proposition~\ref{prop:cy3-on-gamma-not-j}). For a smooth local CY$_3$ $X$
with tilting presentation $(Q, W)$, this yields the full
$\cH(X) \simeq \cH(\Jac(Q, W))$ identification on $\C^3$, the conifold
(Klebanov--Witten $(Q, W)$), local $\bP^2$ (Beilinson cubic), and
$\C^3/\Z_n$ for all $n$.

\subsection*{Why Hall product preservation is forced}

An alternative way to see the transport is purely representation-theoretic:
both $\cH(X)$ and $\cH(A)$ are instances of the \emph{universal positive-
geometry grammar} $Y^+(X) = H^\bullet_{\mathrm{eq}}(\cM^+_{\mathrm{eff}}(X),
\vphi_W)$ (CLAUDE.md, ``Universal positive-geometry grammar''). The
grammar is invariant under any derived equivalence that preserves the
effective semistable moduli stack and the superpotential. Derived Morita
preserves both, by construction. Hence the CoHAs agree, and the Hall product
structure is preserved.

% =======================================================================
\section{Gluing-preservation theorem for mutations}
\label{sec:mutation-preservation}
% =======================================================================

If $T$ and $T'$ are two tilting bundles on the same CY$_3$ $X$, they give two
tilting algebras $A = \End(T)$ and $A' = \End(T')$, and by
Theorem~\ref{thm:hall-product-transport} the same CoHA $\cH(X)$: both
transports land on the same Hall-product object. This induces an isomorphism
$\cH(A) \simeq \cH(A')$. When $T$ and $T'$ are related by a \emph{mutation}
(Derksen--Weyman--Zelevinsky 2008), the induced isomorphism has an explicit
combinatorial presentation.

\subsection*{Mutation at a vertex}

\begin{definition}[Mutation of quiver-with-potential; DWZ 2008,
arXiv:\texttt{0704.0649} Defs.~5.1, 5.5, 7.1]
\label{def:mutation-qp}
\ClaimStatusDefinition\ Let $(Q, W)$ be a reduced quiver-with-potential
(no $2$-cycles in $Q$, no $2$-cycle terms in $W$) with vertex set $I$, and
fix a vertex $k \in I$ supporting no loops. The \emph{mutation at $k$},
$\mu_k(Q, W) = (Q', W')$, factors as a composite of two moves:
the \emph{premutation} $\tilde\mu_k$ followed by the \emph{reduction}
$\mathrm{red}$ of DWZ Theorem~4.6,
\[
  \mu_k(Q, W) \;=\; \mathrm{red}\bigl(\tilde\mu_k(Q, W)\bigr).
\]
The premutation $\tilde\mu_k(Q, W) = (\tilde Q, \tilde W)$ is the three-step
arrow surgery:
\begin{enumerate}
\item[(M1)] \emph{Compose through $k$.} For each pair of arrows
$a \colon i \to k$ and $b \colon k \to j$ with $i, j \neq k$, adjoin a
formal composite $[ba] \colon i \to j$ to the arrow set.
\item[(M2)] \emph{Reverse $k$-incident arrows.} Replace each
$a \colon i \to k$ by $a^{*} \colon k \to i$, and each $b \colon k \to j$
by $b^{*} \colon j \to k$; arrows not incident to $k$ are unchanged.
\item[(M3)] \emph{Update the potential.} Set
\[
  \tilde W \;=\; [W]_{\{[ba]\}} \;+\; \sum_{\substack{a \colon i \to k \\ b \colon k \to j}}
  [ba]\, b^{*}\, a^{*},
\]
where $[W]_{\{[ba]\}}$ is the cyclic representative of $W$ obtained by
replacing every length-two path $ba$ through $k$ by the formal symbol $[ba]$.
\end{enumerate}
The reduction splits the premutation: DWZ Theorem~4.6 provides a canonical
right-equivalence decomposition
$(\tilde Q, \tilde W) \simeq_{\mathrm{r.e.}} (Q_{\mathrm{red}},
W_{\mathrm{red}}) \oplus (Q_{\mathrm{triv}}, W_{\mathrm{triv}})$
into a \emph{reduced} summand (no $2$-cycle terms in $W_{\mathrm{red}}$) and
a \emph{trivial} summand (Jacobi algebra zero), unique up to right-equivalence.
Set $(Q', W') = (Q_{\mathrm{red}}, W_{\mathrm{red}})$.
\end{definition}

The reduction is not the physicist's ``Euler--Lagrange elimination'' but the
Derksen--Weyman--Zelevinsky right-equivalence: a $2$-cycle term
$c\, x\, y$ in $\tilde W$ is removed by the change of variable
$x \mapsto x - c^{-1}\,\partial_{y}(\text{rest of }\tilde W)$, which DWZ
prove is an automorphism of the completed path algebra preserving the cyclic
class of the potential (DWZ Thms.~4.2, 4.6).

The \emph{non-degeneracy hypothesis} on $(Q, W)$ is that every iterated
premutation admits a reduction. DWZ \S7 prove that non-degenerate QPs are
generic and that all QPs arising from local CY$_3$ tilting --- the conifold,
local $\bP^2$, $\C^3/\Z_n$ --- are non-degenerate (DWZ Cor.~7.4).

\begin{proposition}[Mutation is an involution in reduced QPs; DWZ 2008
Thm.~5.7]
\label{prop:mutation-involution}
\ClaimStatusProposition\ In the category whose objects are reduced QPs and
whose morphisms are right-equivalences, $\mu_k \circ \mu_k$ is naturally
right-equivalent to the identity: $\mu_k(\mu_k(Q, W)) \simeq_{\mathrm{r.e.}}
(Q, W)$. The involution holds \emph{up to right-equivalence}, not strictly:
the double premutation $\tilde\mu_k^2$ differs from the identity by a
trivial $2$-cycle summand that cancels only after reduction.
\end{proposition}

Geometrically, $T' = \mu_k(T)$ is \emph{not} a plain ``dual-and-shift''
$T_k^{\vee}[1]$ but a cone on the $\Ext^1$-arrow data at $k$. Replace the
summand $T_k$ by one of two mutation cones,
\[
  T_k^{+} \;=\; \mathrm{cone}\!\left(T_k \;\xrightarrow{\;\mathrm{ev}\;}\;
  \bigoplus_{b \colon k \to j} T_j^{\oplus \dim b}\right),
  \quad
  T_k^{-} \;=\; \mathrm{cone}\!\left(\bigoplus_{a \colon i \to k}
  T_i^{\oplus \dim a} \;\xrightarrow{\;\mathrm{coev}\;}\; T_k\right)[-1],
\]
the first built from arrows \emph{out of} $k$ and the second from arrows
\emph{into} $k$; the two differ by a global shift (Iyama--Reiten 2008,
arXiv:\texttt{0707.1669} \S5; Keller--Yang 2011, arXiv:\texttt{0906.0761}
\S3). The involution Proposition~\ref{prop:mutation-involution} corresponds
geometrically to $\mu_k^{-}(\mu_k^{+}(T)) \simeq T$, which is a
Serre-duality statement on the two cones, \emph{not} to the spurious
$T_k^{\vee}[2] = T_k$.

\begin{theorem}[Mutation gives derived equivalence; Keller--Yang 2011,
arXiv:\texttt{0906.0761} Thms.~3.2 and 6.3]
\label{thm:keller-yang-mutation}
\ClaimStatusTheorem\ Assume $(Q, W)$ is non-degenerate. The mutation
$\mu_k \colon (Q, W) \mapsto (Q', W') = \mu_k(Q, W)$ lifts to a pair of
triangulated equivalences
\[
  \Mu_k^{+} \;:\; \cD(\Gamma(Q, W)) \;\xrightarrow{\;\sim\;}\;
  \cD(\Gamma(Q', W')),
  \qquad
  \Mu_k^{-} \;:\; \cD(\Gamma(Q, W)) \;\xrightarrow{\;\sim\;}\;
  \cD(\Gamma(Q', W')),
\]
where $\cD(\Gamma) = D(\Gamma\text{-}\mathrm{dg\text{-}mod})$. The two equivalences
are constructed from two distinct tilting objects $T_k^{+}, T_k^{-}$ in
$\cD(\Gamma(Q, W))$, built from the decomposition $\Gamma(Q, W) =
\bigoplus_{i \in I} P_i$ into indecomposable projective dg-modules:
\[
  T_k^{+} \;=\; \mathrm{cone}\Bigl(P_k \;\xrightarrow{\;(b)_{b \colon k \to j}\;}\;
  \bigoplus_{b \colon k \to j} P_{t(b)}\Bigr) \;\oplus\; \bigoplus_{i \neq k} P_i,
\]
with the differential given by the arrows $b$ \emph{out of} $k$ in $Q$ (the
sum ranges over $b \in Q_1$ with source $k$ and target $t(b) = j$), and
similarly $T_k^{-}$ built from arrows \emph{into} $k$. Keller--Yang Thm.~3.2
establishes that $T_k^{\pm}$ is a compact generator of $\cD(\Gamma(Q, W))$
with $\Ext^{>0}(T_k^{\pm}, T_k^{\pm}) = 0$; Keller--Yang Thm.~6.3 identifies
$\End_{\cD(\Gamma(Q, W))}(T_k^{\pm}) \simeq \Gamma(Q', W')$ as dg-algebras,
and then Theorem~\ref{thm:derived-morita-tilting} yields the equivalences
$\Mu_k^{\pm}$.

Along $\Mu_k^{\pm}$, the simple $S_k$ (the simple at vertex $k$) over
$\Gamma(Q, W)$ maps to $S_k^{\pm}[\mp 1]$ over $\Gamma(Q', W')$: the
\emph{shift direction is opposite} for $\Mu_k^{+}$ and $\Mu_k^{-}$, and the
two equivalences differ by the spherical twist at $S_k$ (Keller--Yang Prop.~6.4).
\end{theorem}

\subsection*{Induced CoHA automorphism}

\begin{theorem}[Mutation-induced CoHA automorphism]
\label{thm:mutation-coha-automorphism}
\ClaimStatusTheorem\ Let $(Q, W)$ be a non-degenerate quiver-with-potential
whose Ginzburg dg-lift $\Gamma(Q, W)$ is exact $3$-Calabi--Yau, with Jacobi
algebra $A = \Jac(Q, W) = H^0(\Gamma(Q, W))$, and let $\mu_k$ denote mutation
at vertex $k$. Let $\Mu_k^{+} \colon \cD(\Gamma(Q, W)) \to
\cD(\Gamma(\mu_k(Q, W)))$ be the positive Keller--Yang equivalence of
Theorem~\ref{thm:keller-yang-mutation}. Transporting Hall products via
Theorem~\ref{thm:hall-product-transport}:
\[
  \cH(A) \;\xrightarrow{\;\Mu_{k,\star}^{+}\;}\; \cH(A')
\]
is an isomorphism of super-bialgebras. If $(Q, W)$ and $(Q', W') =
\mu_k(Q, W)$ arise as tilting presentations of the \emph{same} geometric
CY$_3$ $X$ --- equivalently, $\mu_k$ relates two tilting bundles $T, T'$ on
$X$ --- then $\Mu_{k,\star}^{+}$ is an automorphism of the single CoHA
$\cH(X)$. The opposite mutation $\Mu_k^{-}$ gives the inverse automorphism
$(\Mu_{k,\star}^{+})^{-1}$.
\end{theorem}

\begin{proof}[Proof.]
By Theorem~\ref{thm:keller-yang-mutation}, $\Mu_k^{+}$ is a triangulated
equivalence between the derived categories of the two CY$_3$ Ginzburg
algebras, constructed as a derived Morita equivalence along the tilting
object $T_k^{+}$. By Theorem~\ref{thm:hall-product-transport}, any derived
Morita equivalence between CY$_3$ Ginzburg algebras induces a super-bialgebra
isomorphism on the critical CoHAs, with equivariant and $\Z/2$-supergrading
structure (DHSM 2023 Thm.~B). When $(Q, W)$ and $(Q', W')$ arise as two
tilting presentations of the same $X$, the corresponding tilting objects
$T, T'$ give two factorisations of the identity
$\mathrm{id}_{\cH(X)} = (\star_{T'})^{-1} \circ \Mu_{k,\star}^{+} \circ \star_T$,
so $\Mu_{k,\star}^{+}$ is an endomorphism of $\cH(X)$. Non-degeneracy of
$(Q, W)$ (DWZ Cor.~7.4) ensures $\Mu_k^{-}$ exists and gives the inverse
at the derived level, hence $\Mu_{k,\star}^{+}$ is an automorphism.
\end{proof}

\subsection*{Geometric interpretation: Seiberg duality}

In the physics literature, mutation at a vertex is \emph{Seiberg duality}: a
quiver gauge theory undergoes an exact IR equivalence under a node mutation
with appropriate matter content adjustment (Seiberg 1994; Berenstein--Douglas
2002). The present theorem is the derived-category / CoHA shadow of this
physical statement: the non-perturbative BPS algebra is invariant under
Seiberg duality.

\begin{example}[Conifold flop via mutation]
\label{ex:conifold-seiberg-duality}
\ClaimStatusDerivation\ The Klebanov--Witten quiver
$(Q_{\mathrm{con}}, W_{\mathrm{con}})$ has two vertices
$I = \{0, 1\}$, four arrows $a_1, a_2 \colon 0 \to 1$ and $b_1, b_2 \colon 1
\to 0$, and the cyclic quartic potential
\[
  W_{\mathrm{con}} \;=\; a_1\, b_1\, a_2\, b_2 \;-\; a_1\, b_2\, a_2\, b_1 \,.
\]
Premutation $\tilde\mu_0$ (DWZ recipe applied at vertex $k = 0$). Step
(M1) requires pairs of arrows $a \colon i \to k$ and $b \colon k \to j$
--- i.e., arrows into $0$ composed with arrows out of $0$. The arrows into
$0$ are $b_1, b_2 \colon 1 \to 0$, and the arrows out of $0$ are $a_1,
a_2 \colon 0 \to 1$; the composites $a \cdot b$ through $0$ run from
$1$ to $1$, adjoined as formal arrows $[a_i b_j] \colon 1 \to 1$ for
$i, j \in \{1, 2\}$ --- four new arrows.
Step (M2) reverses the $0$-incident arrows: $a_i \colon 0 \to 1$ becomes
$a_i^* \colon 1 \to 0$, and $b_j \colon 1 \to 0$ becomes $b_j^* \colon 0
\to 1$. Step (M3) gives
\[
  \tilde W \;=\; [a_1 b_1]\,[a_2 b_2] \;-\; [a_1 b_2]\,[a_2 b_1]
  \;+\; \sum_{i, j}\, [a_i b_j]\, a_i^{*}\, b_j^{*},
\]
where the first two terms are the original $W_{\mathrm{con}}$ rewritten
in the composite variables, and the last sum is the universal $[ba]\, b^*\,
a^*$ piece. The premutation is a cyclic quartic-plus-cubic.

Reduction: the premutation has no $2$-cycles at vertex $0$ or $1$ (the
paths $a_i^* b_j^*$ run $1 \to 0 \to 1$ so sit in the $(1, 1)$-block, but
in pairs $a_i^*, b_j^*$ which are single arrows, not length-two cycles
within a single vertex). However, the cubic terms $[a_i b_j] a_i^* b_j^*$
together with the quartic pair up under the right-equivalence of DWZ \S4.
Explicitly, setting $c_{ij} = [a_i b_j]$, the cyclic derivatives
$\partial_{c_{ij}} \tilde W = \partial_{c_{ij}}\!\bigl(c_{11} c_{22} -
c_{12} c_{21} + \sum_{k\ell} c_{k\ell}\, a_k^* b_\ell^*\bigr)$ solve for
the $c_{ij}$ in terms of $a_i^* b_j^*$, and the reduced potential is
\[
  W' \;=\; a_1^*\, b_1^*\, a_2^*\, b_2^* \;-\; a_1^*\, b_2^*\, a_2^*\, b_1^*,
\]
the same Klebanov--Witten quartic with roles $a \leftrightarrow a^*$ and
$b \leftrightarrow b^*$ swapped. The result $\mu_0(Q_{\mathrm{con}},
W_{\mathrm{con}})$ is right-equivalent to $(Q_{\mathrm{con}}, W_{\mathrm{con}})$
with the vertex labels $\{0, 1\}$ swapped --- the $\Z/2$ symmetry
interchanging the two chambers of the Klebanov--Witten quiver.

The derived equivalence $\Mu_0^{+} \colon \cD(\Gamma(Q_{\mathrm{con}},
W_{\mathrm{con}})) \xrightarrow{\sim} \cD(\Gamma(\mu_0(Q_{\mathrm{con}},
W_{\mathrm{con}})))$ is exactly the Bridgeland--Chen derived equivalence
$D^b(\tilde Y) \simeq D^b(\tilde Y^{+})$ between the two small resolutions
of the conifold singularity $xy = zw$ (Bridgeland 2002, arXiv:\texttt{math/0009053};
Chen 2002). On the Hall-algebra side, $\Mu_{0,\star}^{+} \colon
\cH(\tilde Y) \xrightarrow{\sim} \cH(\tilde Y^{+})$ transports the CoHA
structure; under the identification $\tilde Y \simeq \tilde Y^{+}$ as
abstract varieties (the two small resolutions of the same conifold are
abstractly isomorphic, related by the Atiyah flop), this becomes a
non-trivial automorphism
\[
  \Mu_{0,\star}^{+} \;\in\; \mathrm{Aut}_{\mathrm{bialg}}(\cH(\tilde Y))
\]
equal to the Kontsevich--Soibelman wall-crossing cocycle at the flop wall
(KS 2008 \S6; Nagao 2010, arXiv:\texttt{1002.4884}).

\emph{Braid structure at two vertices.} For $|I| = 2$, the braid group
$B_2 \simeq \Z$ is freely generated by a single element $\mu_0 \mu_1$
(the product of the two mutations, acting as translation on the cluster
tube). There is no ``$\mu_0 \mu_1 \mu_0 = \mu_1 \mu_0 \mu_1$'' relation
because $B_2$ has no braid relation; the correct statement is that at the
conifold, the braid-type subgroup of $\mathrm{Aut}_{\mathrm{bialg}}(\cH(\tilde Y))$
generated by the mutation cocycle is cyclic, and its infinite order
corresponds to the infinite stack of conifold walls in the KS wall-crossing
diagram. This matches Nagao's tilting-theoretic derivation of the conifold
partition function (Nagao--Nakajima, arXiv:\texttt{0809.2992}).
\end{example}

\begin{example}[Local $\bP^2$ nine-arrow Beilinson quiver and $B_3$-braid]
\label{ex:local-p2-mutations}
\ClaimStatusDerivation\ The Beilinson tilting
$T = \cO_{\bP^2} \oplus \cO_{\bP^2}(1) \oplus \cO_{\bP^2}(2)$ on
$K_{\bP^2} = \Tot(\cO_{\bP^2}(-3))$ gives a quiver $Q_{\bP^2}$ on three
vertices $I = \{0, 1, 2\}$, with arrow spaces
\[
  Q_{\bP^2, 1}(i, j) \;=\;
  \begin{cases}
    \Hom(\cO(i), \cO(j)) \;=\; H^{0}(\bP^2, \cO(j - i)) \;=\; \C^{3} &
    \text{if } j = i + 1, \\
    0 & \text{otherwise,}
  \end{cases}
\]
yielding three strands of three parallel arrows $x_\alpha \colon 0 \to 1$,
$y_\alpha \colon 1 \to 2$, and (by descent through the $(-3)$ twist)
$z_\alpha \colon 2 \to 0$, $\alpha = 1, 2, 3$, indexed by a basis of
$H^0(\bP^2, \cO(1)) = \C^3$. The superpotential is the antisymmetric cubic
\[
  W_{\bP^2} \;=\; \sum_{\sigma \in S_3} \mathrm{sgn}(\sigma)\,
  x_{\sigma(1)}\, y_{\sigma(2)}\, z_{\sigma(3)}
\]
(Beilinson 1978; Beasley--Plesser 2001). The Jacobi algebra
$\Jac(Q_{\bP^2}, W_{\bP^2})$ is the NCCR of $\C^3/\Z_3$ in the McKay
convention via Craw--Ishii 2004.

\emph{Three generators.} Mutation at each vertex gives an operator
$\mu_k \in \mathrm{Aut}_{\mathrm{r.e.}}(\text{QP})$ for $k \in \{0, 1, 2\}$.
Because $W_{\bP^2}$ is cyclically symmetric and $Q_{\bP^2}$ is
vertex-transitive under $\Z/3 \subset S_3$, the three $\mu_k$ are pairwise
related by the cyclic automorphism of $Q_{\bP^2}$.

\emph{$B_3$-braid as Bondal--Polishchuk autoequivalences.} The braid group
$B_3 = \langle \mu_0, \mu_1 \mid \mu_0 \mu_1 \mu_0 = \mu_1 \mu_0 \mu_1 \rangle$
(two-generator presentation) acts on $D^b(\mathrm{Coh}\, \bP^2)$ via the
Bondal--Polishchuk exceptional-collection twists $T_{\cO(i)}$
(Bondal--Polishchuk 1993, ``Homological properties of associative algebras:
the method of helices''), extended to $D^b(\mathrm{Coh}\, K_{\bP^2})$ via the
total-space pullback. On the Ginzburg dg-algebra side, each Bondal--Polishchuk
twist $T_i$ corresponds to a spherical twist at the simple $S_i$, and
Keller--Yang 2011 Prop.~6.4 identifies the mutation $\mu_i$ with the inverse
of this spherical twist in $\cD(\Gamma(Q_{\bP^2}, W_{\bP^2}))$.

The braid relation $\mu_0 \mu_1 \mu_0 \simeq_{\mathrm{r.e.}} \mu_1 \mu_0
\mu_1$ is the cluster-combinatorial shadow of the $\mathrm{SL}_2(\Z)$-covariance
of $\bP^2$: the $B_3$-braid is precisely $B_3 = \pi_1(\mathrm{Conf}_3(\C))$,
acting on the $3$-helix $(\cO, \cO(1), \cO(2))$ via the monodromy of the
Kuznetsov exceptional-collection moduli (Kuznetsov 2007, arXiv:\texttt{math/0610957};
Bridgeland 2009, arXiv:\texttt{0909.4299}).

\emph{Transport to CoHA.} By
Theorem~\ref{thm:mutation-coha-automorphism}, each $\mu_k$ transports to an
automorphism $\mu_{k, \star} \in \mathrm{Aut}_{\mathrm{bialg}}(\cH(K_{\bP^2}))$,
and the braid relation $\mu_{0,\star} \mu_{1,\star} \mu_{0,\star} =
\mu_{1,\star} \mu_{0,\star} \mu_{1,\star}$ holds in the automorphism group.
Under the identification $\cH(K_{\bP^2}) \simeq Y_{1, 0, 0}[\psi]$ with the
Gaiotto--Rap\v{c}\'ak corner VOA (Gaiotto--Rap\v{c}\'ak 2018,
arXiv:\texttt{1703.00982}), this $B_3$-braid realises the triality
$\epsilon_1 \leftrightarrow \epsilon_2 \leftrightarrow \epsilon_3$ of the
three toric $\C^\times$-weights at the McKay vertex of $\C^3/\Z_3$.
\end{example}

\subsection*{The mutation cocycle as a braid-group representation}

The mutations $\{\mu_k\}_{k \in I}$ satisfy the braid relations of $B_{|I|}$
at two distinct layers: at the level of derived equivalences, via Keller--Yang
Theorem~\ref{thm:keller-yang-mutation}; and at the level of CoHA bialgebra
automorphisms, via transport through
Theorem~\ref{thm:mutation-coha-automorphism}. The combined statement is a
bialgebra representation of the braid group.

\begin{theorem}[Mutation cocycle; braid-group representation on the CoHA]
\label{thm:mutation-cocycle}
\ClaimStatusTheorem\ Let $X$ be a smooth local CY$_3$ with tilting bundle
$T = \bigoplus_{i \in I} T_i$ and non-degenerate QP presentation $(Q, W)$.
Denote by $B_{|I|}$ the Artin braid group on $|I|$ strands, with standard
generators $\sigma_0, \ldots, \sigma_{|I| - 2}$ subject to
\[
  \sigma_i \sigma_{i+1} \sigma_i \;=\; \sigma_{i+1} \sigma_i \sigma_{i+1},
  \qquad
  \sigma_i \sigma_j \;=\; \sigma_j \sigma_i \quad (|i - j| \geq 2).
\]
The assignment $\sigma_k \mapsto \Mu_k^{+}$ (Keller--Yang equivalence at
vertex $k$) extends to a group homomorphism
\[
  \rho_X^{\mathrm{der}} \;:\; B_{|I|} \;\longrightarrow\;
  \mathrm{Auteq}\bigl(\cD(\Gamma(Q, W))\bigr)
\]
into the group of triangulated autoequivalences, and, composing with the
Hall-product transport $\Mu_{k,\star}$ of
Theorem~\ref{thm:mutation-coha-automorphism}, to
\[
  \rho_X \;=\; (-)_{\star} \circ \rho_X^{\mathrm{der}} \;:\; B_{|I|}
  \;\longrightarrow\; \mathrm{Aut}_{\mathrm{bialg}}(\cH(X)),
  \qquad \sigma_k \;\longmapsto\; \Mu_{k,\star}^{+}.
\]
The image $\rho_X(B_{|I|})$ is the \emph{mutation cocycle subgroup} of
bialgebra automorphisms of $\cH(X)$.
\end{theorem}

\begin{proof}[Proof.]
Two parts: the braid relations at the derived level, and the functoriality
of Hall-product transport.

(i) \emph{Braid relations at the derived level.} By
Theorem~\ref{thm:keller-yang-mutation}, each $\Mu_k^{+}$ is the derived
equivalence induced by the tilting object $T_k^{+} = \mathrm{cone}(P_k \to
\bigoplus_{b \colon k \to j} P_{t(b)}) \oplus \bigoplus_{i \neq k} P_i$.
Keller--Yang 2011 Prop.~6.4 identifies $\Mu_k^{+}$ with the inverse of the
spherical twist $\mathrm{Tw}_{S_k}^{-1}$ at the simple module $S_k$ over
$\Gamma(Q, W)$. The spherical-twist functors $\mathrm{Tw}_{S_k}$ satisfy
the Seidel--Thomas braid relations
\[
  \mathrm{Tw}_{S_k} \mathrm{Tw}_{S_\ell} \mathrm{Tw}_{S_k}
  \;=\;
  \mathrm{Tw}_{S_\ell} \mathrm{Tw}_{S_k} \mathrm{Tw}_{S_\ell}
  \quad\text{if}\quad \mathrm{Ext}^{*}(S_k, S_\ell) = \C[-1] \oplus \C[-2],
\]
i.e., when the two simples are ``adjacent'' in the sense of Seidel--Thomas
2001 arXiv:\texttt{math/0001043} Thm.~1.2, and commute when
$\mathrm{Ext}^{*}(S_k, S_\ell) = 0$. For $\Gamma(Q, W)$ a CY$_3$ Ginzburg
algebra, Keller--Yang Thm.~6.4 verifies that $\Ext^{*}(S_k, S_\ell)$ has
this form precisely when $k$ and $\ell$ are connected by an arrow in $Q$
(i.e.\ adjacent in the $2$-coloured exchange graph); the commutation holds
when $k, \ell$ are disconnected.

Therefore the assignment $\sigma_k \mapsto \Mu_k^{+}$ respects both the
braid relation (for adjacent $k, \ell$) and the commutation relation (for
non-adjacent), and extends to the group homomorphism $\rho_X^{\mathrm{der}}$
on $B_{|I|}$.

(ii) \emph{Functoriality of Hall-product transport.} By
Theorem~\ref{thm:hall-product-transport}, every derived Morita equivalence
between two CY$_3$ Ginzburg algebras induces a bialgebra isomorphism on the
critical CoHAs. This assignment is functorial: the composition
$\Mu_k^{+} \circ \Mu_\ell^{+}$ transports to $\Mu_{k,\star}^{+} \circ
\Mu_{\ell,\star}^{+}$ by the universal property of critical cohomology with
respect to derived Morita, which descends from the functoriality of
DHSM 2023 arXiv:\texttt{2303.12592} Thm.~A.

Combining (i) and (ii): $\rho_X = (-)_{\star} \circ \rho_X^{\mathrm{der}}$
is a group homomorphism from $B_{|I|}$ to
$\mathrm{Aut}_{\mathrm{bialg}}(\cH(X))$.
\end{proof}

\begin{corollary}[Chamber-independence of the CoHA bialgebra]
\label{cor:chamber-independence}
\ClaimStatusTheorem\ Let $X$ be as above. The CoHA bialgebra $\cH(X)$,
viewed as an abstract bialgebra, is \emph{invariant} under the mutation
cocycle: for any $\gamma \in B_{|I|}$ and any homogeneous elements $x, y
\in \cH(X)$, writing $\star$ for the Hall product,
\[
  \rho_X(\gamma)(x \star y) \;=\; \rho_X(\gamma)(x) \star \rho_X(\gamma)(y),
  \qquad
  \rho_X(\gamma) \Delta \;=\; \Delta\, \rho_X(\gamma).
\]
The chamber-dependent structure carried by $\cH(X)$ is not the bialgebra
itself but the \emph{PBW/BPS decomposition} of $\cH(X)$: in each chamber
$\mathfrak{C} \subset \mathrm{Stab}(X)$, the BPS sheaf
$\mathcal{BPS}_{\mathfrak{C}}$ on $\cM(X)$ picks out a different
$t$-structure heart, and the PBW basis of $\cH(X)$ depends on this choice.
The mutation cocycle translates between PBW bases but leaves the bialgebra
structure constants invariant.
\end{corollary}

\begin{proof}[Proof.]
Invariance under each generator $\sigma_k$: $\Mu_{k,\star}^{+}$ is by
Theorem~\ref{thm:mutation-coha-automorphism} a bialgebra isomorphism. The
braid-group extension $\rho_X$ is a composition of generators, hence
bialgebra-preserving.

Chamber dependence of the PBW basis: KS 2008 \S6 proves that the BPS sheaf
$\mathcal{BPS}_{\mathfrak{C}}$ depends on the stability condition
$\mathfrak{C}$, and that the PBW-type basis of $\cH(X)$ indexed by
BPS states is chamber-indexed. Davison--Meinhardt 2020
arXiv:\texttt{1601.02479} Thm.~A establishes that the \emph{bialgebra}
$\cH(X)$ is chamber-independent, with the PBW decomposition providing the
chamber-dependent structure.
\end{proof}

The mutation cocycle is the algebraic shadow of the Donaldson--Thomas
wall-crossing formula (Kontsevich--Soibelman 2008 \S6): each wall-crossing
transformation is a chamber-to-chamber BPS reindexing, and the mutation
cocycle realises this at the level of bialgebra automorphisms, preserving
the underlying $\cH(X)$.

% =======================================================================
\section*{Summary and placement}
% =======================================================================

Part IV has collapsed the \v{C}ech and McKay atlases of Parts I--III onto a
single finite-dimensional algebra $A = \End(T)^{\mathrm{op}}$ with Ginzburg
dg-lift $\Gamma(Q, W)$ via the Van den Bergh derived Morita equivalence
(Theorem~\ref{thm:vdb-equivalence}). The CY$_3$ property is carried by the
dg-lift $\Gamma$, not by its degree-zero cohomology $\Jac(Q, W) = H^0(\Gamma)$
(Proposition~\ref{prop:cy3-on-gamma-not-j}). The tilting bundle
$T = \bigoplus_i T_i$ serves as universal gluing data
(Theorem~\ref{thm:tilting-jacobi}): all chart-wise transitions fold into the
Jacobi relations $\partial_a W = 0$. The derived Morita equivalence is the
adjoint pair $L \dashv R$ satisfying compact generation and higher-Ext
vanishing (hypotheses (H1)--(H2) of Theorem~\ref{thm:derived-morita-tilting}).
The Hall product transports along any such equivalence between tilting algebras
(Theorems~\ref{thm:moduli-morita}--\ref{thm:hall-product-transport}, DHSM 2023
arXiv:\texttt{2303.12592} Thms.~A--B), and the
Derksen--Weyman--Zelevinsky--Keller--Yang mutation machinery
(Theorems~\ref{thm:keller-yang-mutation}--\ref{thm:mutation-cocycle},
arXiv:\texttt{0704.0649} and arXiv:\texttt{0906.0761})
realises the Artin braid group $B_{|I|}$ as explicit CoHA bialgebra
automorphisms $\rho_X \colon B_{|I|} \to \mathrm{Aut}_{\mathrm{bialg}}(\cH(X))$,
with chamber-independence
(Corollary~\ref{cor:chamber-independence}) of the bialgebra structure under
this action. The
four subsections above furnish the tilting toolkit that downstream parts of
the gluing chapter deploy: Part V (factorisation-algebra gluing) extends the
tilting picture to non-tilting-admitting varieties via local factorisation
presentations; Part VI (wall-crossing) frames the mutation cocycle as the
chamber-to-chamber KS quantum-dilogarithm factorisation; Part VIII (the
K3 $\times$ E master example) invokes the \emph{Serre-equivariant quasi-NCCR}
as the obstruction-aware analogue of Van den Bergh, where a global tilting
fails but a locally-defined Serre-equivariant tilting patches via the
factorisation pushforward.
